\documentclass[11pt,a4paper]{article}

\usepackage[T1]{fontenc}
\usepackage[utf8]{inputenc}
\usepackage{lmodern}
\usepackage[a4paper,margin=26mm]{geometry}
\usepackage{amsmath,amssymb,amsthm,mathtools}
\usepackage{booktabs,array}
\usepackage{microtype}
\usepackage{xcolor}
\usepackage{enumitem}
\usepackage{url}
\usepackage[colorlinks=true,linkcolor=blue!55!black,citecolor=blue!55!black,
            urlcolor=blue!55!black]{hyperref}

\newtheorem{theorem}{Theorem}
\newtheorem{lemma}[theorem]{Lemma}
\newtheorem{proposition}[theorem]{Proposition}
\newtheorem{remark}[theorem]{Remark}
\newcommand{\F}{\mathbb F}
\newcommand{\bits}{\{0,1\}}
\newcommand{\xor}{\mathbin{\oplus}}
\newcommand{\concat}{\mathbin\Vert}
\newcommand{\Sign}{\ensuremath{\mathsf{Sign}}}
\newcommand{\Verify}{\ensuremath{\mathsf{Verify}}}
\newcommand{\KeyGen}{\ensuremath{\mathsf{KeyGen}}}
\newcommand{\pk}{\ensuremath{\mathsf{pk}}}
\newcommand{\sk}{\ensuremath{\mathsf{sk}}}
\newcommand{\MsgHash}{\ensuremath{\mathsf{MsgHash}}}
\newcommand{\CoinHash}{\ensuremath{\mathsf{CoinHash}}}
\newcommand{\TCCR}{\ensuremath{\mathsf{TCCR}}}
\newcommand{\EncodeNoise}{\ensuremath{\mathsf{EncodeNoise}}}
\newcommand{\Prb}{\mathop{\rm Pr}}
\newcommand{\sample}{\mathrel{\leftarrow\$}}

\title{\bfseries Cross-Profile Seed-Tree Reuse in ReSolveD-$\alpha$:\\
Recovery of an Equivalent Signing Witness and Fresh-Message Forgery\\
from Same-Key, Same-Message Deterministic S/F Signatures\\
at One Security Level}
\author{Martin Feussner\\
\small Selmer Center, University of Bergen\\
\small \texttt{martin.feussner@uib.no}}
\date{1 October 2026}

\begin{document}
\maketitle

\begin{center}
\fbox{\parbox{0.92\linewidth}{\small\textbf{Provenance and review status.}
The ReSolveD-$\alpha$ attack, derivation, implementation, experiments, and
initial manuscript were produced by OpenAI Codex using Daybreak Blue at Ultra
reasoning effort during an AI-run audit of NGCC signature candidates.  The
mechanism was found after the same audit had identified a related cross-variant
failure in Galas; no human first proposed its application to ReSolveD-$\alpha$.
A separately instructed hostile-review agent rebuilt the recovery, tried to
disprove the result, and found that the conditional claim survived that review.
At the date above, no human cryptanalyst has independently verified the
argument, implementation, measurements, or novelty assessment.  This report is
a preliminary disclosure for human review and reproduction.}}
\end{center}
\vspace{0.7em}

\begin{abstract}
ReSolveD-$\alpha$ supplies short S and fast F profiles at each of four security
parameters.  At a fixed level they use the same RSD public relation, the same
untagged key formats, and a deterministic signing mode expressly defined by
$\rho=0^\lambda$.  We show a practical cross-profile composition attack
enabled by absent cryptographic separation of S/F keys and derivation domains
under same-key reuse.  At one fixed level, it applies when the same full secret
key---and hence the same effective RSD witness material---signs the same message
once under S and once under F with equal $\rho$.  The deterministic mode supplies
that equality.  The specification inserts no profile identifier in either
\[
 (\mu,s)=\MsgHash(\pk,M),\qquad
 (r,iv)=\CoinHash(\sk,\mu,\rho),
\]
so the two profiles begin from the same root seed, tree salt, and IV.

The F seed tree is much smaller.  Every parent of an F leaf has the same seed as
the S node at the corresponding heap index, although the F and S leaf seeds
themselves generally differ because the profiles use different expansion rules
there.  If the S opening permits reconstruction of the parent of one F-hidden
leaf, the public F bottom rule recovers that leaf.  Completing any one F VOLE
instance gives $u_a$; for $a=0$ this is $u_0$, while for $a>0$ the serialized
correction gives $u_0=u_a\xor c_a$.  The signature then exposes the compact RSD
witness
\[
       W=d\xor u_0[0..\ell_{\rm wit}).
\]
This $W$ is an equivalent signing witness and a complete signing capability:
decoding it yields a regular error $e$ with $He=y$, so an attacker can generate
an honest proof on a fresh message without recovering the original secret seed.

End-to-end experiments recovered the exact honest witness and produced
fresh-message S signatures accepted by verifiers compiled from untouched
submitted verification sources at 160, 256, 384, and 512 bits.  The lead
harness's complete fixed runs took 8.69--102.10 seconds and at most 34 MiB on an
AMD EPYC host; the variation is dominated by transcript grinding.  A separate public-input-only
160-bit implementation recovered in 3.99 seconds and forged in 1.92 seconds.
A 256-key/message real-signer sweep recovered and forged in 256/256 cases, and
all 256 fresh forgeries were accepted by the untouched verifier.  Wrong-message,
wrong-key, different-message, different-$\rho$, F-only, and mutated-witness
controls failed as expected.

The result requires the same security level, the same full secret key and
effective RSD witness material, the same message, and equal $\rho$ across S and
F.  The deterministic mode makes equality of $\rho$ automatic.  The
specification does not instruct users to reuse S/F keys, but it neither
cryptographically prevents that reuse nor separates the S/F key namespaces.
This report does not claim an isolated single-profile, single-oracle EUF-CMA
break or an attack on independently typed S/F keys.
\end{abstract}

\section{Introduction}

ReSolveD-$\alpha$ is a code-based Fiat--Shamir signature that proves knowledge
of a regular syndrome-decoding (RSD) witness using VOLE-in-the-Head,
QuickSilver, and a batch all-but-one vector commitment (BAVC)
\cite{resolved-spec,resolved-paper}.  Each recommended security level has an S
profile with shorter signatures and an F profile intended for faster signing.
The profiles retain one RSD relation while changing the number and sizes of the
VOLE seed trees.

That performance tradeoff creates two different views of the same deterministic
seed tree.  At one fixed level, equal full-key material, message, and $\rho$
give equal roots because the key and coin derivations omit the S/F label.  The
small F tree switches to its
bottom two-hash expansion while the much larger S tree is still in its upper
correlation-robust expansion.  An F leaf is therefore not itself an S leaf with
the same value.  Its parent, however, is a shared internal node.  An S opening
that reconstructs that parent completes the one leaf hidden by the F opening.

The consequence is stronger than transcript malleability.  One completed F
VOLE instance reveals the mask applied to the RSD witness.  The recovered
witness is publicly checkable and is sufficient to create fresh-message
signatures.  No RSD decoding, hash inversion, Fiat--Shamir fixed point, or
weakness in the symmetric primitives is required.

Our contributions are as follows.
\begin{enumerate}[leftmargin=*,itemsep=2pt]
  \item We prove the exact common-node identity and check the required heap
        bounds for every recommended S/F pair.
  \item We give a public recovery algorithm that converts one complementary
        S/F pair into an equivalent RSD signing witness and validates it using
        only the public key.
  \item We reproduce exact witness recovery and fresh-message forgery at all
        four levels, with final acceptance by untouched submitted verifiers and
        adversarial negative controls.
  \item We separate an exact unconditioned index calculation from both the
        opening-conditioned geometry and a 256-key/message real-signer sweep.
  \item We delimit the result as a same-level cross-profile attack requiring
        the same full key/effective witness material, message, and $\rho$.  It
        is not a claim against an isolated S or F signing oracle.
\end{enumerate}

The finding is a practical cross-profile composition attack enabled by absent
cryptographic separation of S/F keys and derivation domains under same-key
reuse.  The
specification neither instructs cross-profile key reuse nor cryptographically
prevents it or separates the S/F namespaces.  The submitted signer omits part
of the specified \CoinHash{} input; we repaired that deviation before the main
experiments, and the attack survives the repair.  A separate artifact-risk
control deliberately reset the untouched S/F APIs' DRNG state to induce equal
nonzero $\rho$.  That test shows that the attack is not introduced by the
repaired wrapper; it is not a claim about normal randomized API use.

\section{Relevant construction}

\subsection{Keys and the equivalent witness relation}

For a fixed level, let $H=[I_{m-k}\mid H_B]\in\F_2^{(m-k)\times m}$.
The RSD block size is six and $w=m/6$.  Key generation samples two independent
$\lambda$-bit strings and computes
\begin{align}
 \mathsf{seed}_{pk},\mathsf{seed}_{sk}&\sample\bits^\lambda,
 &H_B&=\mathsf{SampleMatrix}(\mathsf{seed}_{pk}),\nonumber\\
 e&=\mathsf{SampleNoise}(\mathsf{seed}_{sk}),
 &y&=He,                                                     \label{eq:keygen}\\
 \pk&=(\mathsf{seed}_{pk},y),
 &\sk&=(\mathsf{seed}_{pk},\mathsf{seed}_{sk}).\nonumber
\end{align}
Every consecutive six-bit block of $e$ has Hamming weight one.

Write $e=(e_A,e_B)$ according to the systematic split.  The compact witness is
$W=\EncodeNoise(e)$.  It stores the first five bits of each six-bit block of
$e_B$, and the omitted bit restores odd parity.  Given $(\pk,W)$, anyone can
reconstruct $e_B$ and compute
\begin{equation}
                  e_A=y\xor H_Be_B.                         \label{eq:decode}
\end{equation}
Consequently, any $W$ for which every reconstructed block of $e$ has weight one
is an equivalent signing witness.  It need not lie in the image of
$\mathsf{SampleNoise}$, and recovering $\mathsf{seed}_{sk}$ is unnecessary.

The four recommended RSD tuples and profile sizes are shown in
Table~\ref{tab:params}.  At one level, S and F share $(m,k,w)$, the serialized
key format, and the public witness relation.  Independent \KeyGen{} calls still
sample independent keys; the attack assumes reuse of the same raw full
secret-key material across the two profiles.

\begin{table}[ht]
\centering
\caption{Recommended ReSolveD-$\alpha$ parameters used by the attack.}
\label{tab:params}
\begin{tabular}{lrrrrrrrr}
\toprule
$\lambda$ & $(m,k,w)$ & $\tau_S$ & $\tau_F$ & $L_S$ & $L_F$
 & $|W|$ bytes & $|\sigma_S|$ & $|\sigma_F|$\\
\midrule
160 & $(1860,1056,310)$ & 14 & 21 & 28,672 & 3,328 & 110 & 5,307 & 6,851\\
256 & $(2988,1698,498)$ & 22 & 35 & 49,152 & 4,864 & 177 & 13,973 & 17,932\\
384 & $(4452,2532,742)$ & 34 & 53 & 69,632 & 7,296 & 264 & 31,575 & 40,469\\
512 & $(5940,3378,990)$ & 46 & 72 & 94,208 & 9,216 & 352 & 56,239 & 72,611\\
\bottomrule
\end{tabular}
\end{table}

\subsection{Signing state and witness masking}

For message $M$, the specified signer derives
\begin{align}
 (\mu,s)&=\mathsf{XOF}(0x00\concat\mathsf{seed}_{pk}\concat y\concat M;3\lambda),
                                                               \label{eq:msghash}\\
 (r,iv)&=\mathsf{XOF}(0x01\concat\sk\concat\mu\concat\rho;
                      \lambda+120).                            \label{eq:coinhash}
\end{align}
The signer normally samples $\rho\sample\bits^\lambda$; Section~4.2.2 of the
submission expressly defines deterministic signing by setting
$\rho=0^\lambda$.  Neither equation contains the profile identifier.

The profile-specific BAVC and VOLE algorithms expand $(r,s,iv)$ into instances
$u_0,\ldots,u_{\tau-1}$.  The signature serializes corrections
\begin{equation}
                         c_a=u_0\xor u_a\qquad(a\geq1)          \label{eq:corr}
\end{equation}
and masks the RSD witness as
\begin{equation}
                         d=W\xor u_0[0..\ell_{\rm wit}).       \label{eq:mask}
\end{equation}
The BAVC opening reveals every seed except one leaf per VOLE instance, while
the hidden leaf commitments remain serialized.  Fiat--Shamir hashes bind the
commitment, corrections, $d$, and QuickSilver proof into the final challenge.

Verification reconstructs the opened VOLE view, recomputes all three
Fiat--Shamir challenges, and checks the QuickSilver proof for the public RSD
relation.  The proof coins are not recomputed from $\mathsf{seed}_{sk}$.
Therefore an attacker who knows any valid $W$ can execute the public prover
with attacker-chosen coins and make a conventional fresh-message signature.

\subsection{The pcGGM heap}

A profile with $L$ physical leaves uses a one-based heap with root $R_1$,
internal nodes $R_1,\ldots,R_{L-1}$, and leaves
$R_L,\ldots,R_{2L-1}$.  For public $(s,iv)$ its upper expansion is
\begin{align}
 R_{2i}&=\TCCR(R_i,s,iv),\nonumber\\
 R_{2i+1}&=R_{2i}\xor R_i, && i<L/2,                           \label{eq:upper}
\end{align}
and its bottom expansion is
\begin{align}
 R_{2i}&=\TCCR(R_i\xor\mathsf{const}_0,s,iv),\nonumber\\
 R_{2i+1}&=\TCCR(R_i\xor\mathsf{const}_1,s,iv),
                       && L/2\leq i<L.                         \label{eq:bottom}
\end{align}
The final challenge selects one hidden leaf in every logical VOLE instance.
The opening serializes a minimum span of node seeds that covers all other
leaves.  It therefore permits reconstruction of precisely the nodes outside
the union of paths from the hidden leaves to the root.

\section{The cross-profile identity}

\begin{lemma}[Shared deterministic state]\label{lem:state}
Fix one level, a key pair $(\pk,\sk)$, and a message $M$.  S and F signing with
the same $\rho$ produce identical $(\mu,s,r,iv)$.  In particular, this holds
for the specified deterministic mode $\rho=0^\lambda$.
\end{lemma}
\begin{proof}
At a fixed level the profiles have byte-identical key formats and public RSD
parameters.  Equations~\eqref{eq:msghash} and~\eqref{eq:coinhash} are
deterministic functions of the displayed values and contain no S/F identifier.
Their inputs and outputs are therefore equal.
\end{proof}

\begin{lemma}[Shared F internal nodes]\label{lem:internal}
Under Lemma~\ref{lem:state},
\[
                    R_i^F=R_i^S\qquad(1\leq i<L_F).
\]
An F leaf seed is generally not equal to the S seed at the same heap index.
\end{lemma}
\begin{proof}
The roots agree.  Suppose $1<i<L_F$.  The parent that generates $R_i$ has
index strictly below $L_F/2$.  Both profiles therefore apply the upper rule
in~\eqref{eq:upper} at that parent, with the same seed, $s$, and $iv$.
Induction proves equality of every F internal node.

For an F leaf, its parent lies at or below the point where F switches to the
bottom rule~\eqref{eq:bottom}, while S is still using the upper rule.  The two
child computations differ, so equality of the leaf seeds does not follow.
\end{proof}

\begin{proposition}[All F leaf parents are usable S nodes]
\label{prop:bounds}
For every recommended pair, every F leaf position is an internal S heap
position, and the seed of its parent is shared by Lemma~\ref{lem:internal}.
\end{proposition}
\begin{proof}
In zero-based notation the F leaves occupy positions
$L_F-1,\ldots,2L_F-2$, while the last S internal position is $L_S-2$.
Table~\ref{tab:bounds} checks the stronger comparison requested directly from
the submitted schedules:
\[
                       2L_F-2<L_S-1.
\]
Moreover, an F leaf parent has one-based index in
$[L_F/2,L_F-1]$.  Its value is an F internal seed and hence shared by
Lemma~\ref{lem:internal}; each listed $L_F-1<L_S/2$, so S has not reached its
bottom rule at that parent.
\end{proof}

\begin{table}[ht]
\centering
\caption{Self-contained heap bounds for all recommended profile pairs.}
\label{tab:bounds}
\begin{tabular}{lrrrr}
\toprule
Pair & $2L_F-2$ & $L_S-1$ & $L_F-1$ & $L_S/2$\\
\midrule
160S/F & 6,654 & 28,671 & 3,327 & 14,336\\
256S/F & 9,726 & 49,151 & 4,863 & 24,576\\
384S/F & 14,590 & 69,631 & 7,295 & 34,816\\
512S/F & 18,430 & 94,207 & 9,215 & 47,104\\
\bottomrule
\end{tabular}
\end{table}

The distinction in Lemma~\ref{lem:internal} is important.  The attack does not
reinterpret an S node as an equal F leaf.  It asks the S opening to reconstruct
the shared \emph{parent} seed and then applies the public F bottom rule to that
seed.

\section{Equivalent signing witness recovery and forgery}

Let $h_a$ be the F-hidden physical leaf for logical instance $a$, and let
$p_a=\lfloor h_a/2\rfloor$.  Decode both final challenges and reconstruct the
nodes permitted by each opening.  If $p_a$ is outside the S hidden-path union,
the S opening permits reconstruction of $R_{p_a}^S$.  By
Lemma~\ref{lem:internal}, this is $R_{p_a}^F$.  Applying the left or right branch
of~\eqref{eq:bottom} gives the missing F leaf seed.

The F opening already reconstructs every other leaf seed in instance $a$.
Public PRG expansion and XOR therefore give the complete $u_a$.  Write $N_a$
for the number of F leaf seeds contributing to that instance,
$\widehat{\ell}$ for the bit length expanded from each seed, and $b$ for the
number of output bits charged to one underlying PRG primitive block.  Equation
\eqref{eq:corr} gives
\begin{equation}
 u_0=\begin{cases}
       u_a,&a=0,\\
       u_a\xor c_a,&a>0,
     \end{cases}                                             \label{eq:u0}
\end{equation}
and the witness follows from Equation~\eqref{eq:mask}:
\begin{equation}
                        W=d\xor u_0[0..\ell_{\rm wit}).       \label{eq:recover}
\end{equation}

\begin{theorem}[Conditional recovery of an equivalent signing witness]
\label{thm:attack}
Suppose, at one fixed security level, the same ReSolveD-$\alpha$ full secret
key---and hence the same effective RSD witness material---signs one message
under S and F with equal $\rho$.  If at least one F hidden-leaf parent lies
outside the S hidden-path union, the two public signatures determine an
equivalent RSD signing witness in
$O(L_S+L_F+N_a\lceil\widehat{\ell}/b\rceil)$ primitive-block work and
$O((L_S+L_F)\lambda)$ bits of memory.
\end{theorem}
\begin{proof}
The equal signing state follows from Lemma~\ref{lem:state}.  For the available
instance, Proposition~\ref{prop:bounds} and the S opening reconstruct the shared
parent, and the F rule reconstructs the missing leaf.  All inputs to the PRG
expansions, corrections, and Equations~\eqref{eq:u0}--\eqref{eq:recover} are
public, so they recover $W$ without secret state.  Decoding $W$ with
Equation~\eqref{eq:decode} and checking each weight-one block validates the
public relation.  Traversing the trees costs $O(L_S+L_F)$ primitive work, and
expanding the $N_a$ F leaf seeds costs
$O(N_a\lceil\widehat{\ell}/b\rceil)$ primitive blocks.  Streaming those leaf
images while storing the two seed heaps uses
$O((L_S+L_F)\lambda)$ bits.
\end{proof}

The complete chain is
\[
\boxed{\begin{array}{c}
\text{missing S/F separation}\Longrightarrow\text{equal roots and shared parents}\\
\Longrightarrow\text{complementary opening}\Longrightarrow u_a
\Longrightarrow u_0\Longrightarrow W\Longrightarrow\text{fresh-message proof.}
\end{array}}
\]

The public validation step is material.  The attacker reconstructs $e$ from
$(\pk,W)$, checks every six-bit block has weight one, and checks $He=y$.  A
successful pair is therefore recognizable without the original
$\mathsf{seed}_{sk}$.  To forge, the attacker runs the ordinary prover on a
fresh message using $W$ and attacker-chosen proof coins.  The verifier checks
the public relation and has no way to require that $W$ came from the original
noise-sampling seed.

\section{Success probability and retry strategy}

For a fixed F parent $p$, let
\[
 c_b(p)=\#\{j:p\text{ is an ancestor of the S physical leaf for }(b,j)\}.
\]
If the decoded S indices are independent and uniform before grinding and
opening-cap rejection, the exact probability that the S opening reconstructs
$p$ is
\begin{equation}
 P_p=\prod_{b=0}^{\tau_S-1}
       \left(1-\frac{c_b(p)}{N_{S,b}}\right).                 \label{eq:pp}
\end{equation}
Averaging Equation~\eqref{eq:pp} over the uniform hidden index of F instance
zero gives the second column of Table~\ref{tab:prob}.

\begin{table}[ht]
\centering
\caption{Analytic and empirical opening geometry.  The analytic column is an
unconditioned independent-index model, not an exact serialized-signature
probability.}
\label{tab:prob}
\begin{tabular}{lrrr}
\toprule
Pair & analytic $\Pr[p_0\text{ exposed}]$
     & conditioned $p_0$ & conditioned any instance\\
\midrule
160S/F & 0.991823130948 & 992/1,000 & 1,000/1,000\\
256S/F & 0.991696001787 & 988/1,000 & 1,000/1,000\\
384S/F & 0.991401471498 & 987/1,000 & 1,000/1,000\\
512S/F & 0.990043848760 & 990/1,000 & 1,000/1,000\\
\bottomrule
\end{tabular}
\end{table}

The qualification is essential.  A real signer serializes only challenges
whose proof-of-work suffix is zero and whose minimum BAVC opening fits
$T_{\rm open}$.  Those conditions alter the index distribution, and S/F
challenges are generated from related transcripts.  The last two columns come
from 1,000 independently sampled challenge pairs per level after applying the
exact submitted opening-size condition.  They are empirical conditioned counts,
not proofs of the real-signer probability.

The real event is publicly testable: attempt recovery for each F instance and
validate the resulting $W$ against the public RSD relation.  If no parent is
available, request another deterministic S/F pair on a new common message.
Changing the message changes $\mu$ and the root except with negligible hash
collision probability.

We also tested the complete signer rather than only its decoded-index geometry.
Across 256 distinct deterministic test keys and messages at 160 bits, recovery
and fresh-message forgery succeeded in 256/256 cases.  Instance zero was enough
in 255 cases; the remaining case succeeded with instance one.  This experiment
supports practical retry behavior, but it does not turn the unconditioned model
into an exact probability theorem.

\section{End-to-end experiments}

\subsection{Specification repairs and verifier separation}

The official specification, not accidental source behavior, defines the design.
The main experiment began from the submitted reference implementation and made
one signer conformance repair: the source's \texttt{hash\_r\_iv} absorbed only
$\mathsf{seed}_{sk}$, whereas Equation~\eqref{eq:coinhash} specifies the full
$\sk=\mathsf{seed}_{pk}\concat\mathsf{seed}_{sk}$.  The repaired signer inserts
$\mathsf{seed}_{pk}$ immediately before $\mathsf{seed}_{sk}$.  It does not add
an S/F identifier and therefore preserves the cross-profile equality.

Two additional helper functions expose public operations already performed by
verification: reconstructing a permitted BAVC node and XORing the PRG images of
a completed VOLE instance.  They receive only serialized signatures, public
parameters, public keys, and messages.  They do not read the secret key, honest
witness, root, or retained signer state.

The submission also contains non-executable or contradictory notation for the
QuickSilver packing, VOLE row length, and byte padding.  The experiment uses the
source-consistent intended transcript and records those interpretations in the
reproducer.  The attack does not exploit them: the recovered $W$ equals the
honest regular witness, satisfies the public RSD relation directly, and can be
used with either the submitted constraint implementation or the natural repaired
specification.  The extraction itself stops before QuickSilver.

Final acceptance was tested by a separate executable linked directly against
untouched submitted verification sources.  No conformance-repaired or
instrumented source file is linked into that binary.  Because \CoinHash{} is a
signer-only coin derivation and is never recomputed by verification, its repair
cannot cause a false verifier acceptance.

\subsection{All-level recovery and forgery}

For each recommended level, the experiment generated one S and one F signature
under the same key and message with the specified deterministic
$\rho=0^\lambda$.  It verified both signatures, recovered $W$, compared it to
the independently retained honest witness, and signed a fresh message from the
recovered witness with attacker-selected proof coins.  The untouched submitted
S verifier accepted every fresh signature and rejected that signature under the
original query message.

Table~\ref{tab:results} reports complete process measurements, including the
two honest query signatures, recovery, fresh proof generation, and controls.
Signing grinds both a challenge suffix and an opening-size condition, so fixed
inputs vary substantially; the 512-bit case happened to accept unusually
quickly and should not be read as a stable benchmark.

\begin{table}[ht]
\centering
\caption{Lead-harness, specification-repaired end-to-end experiments on the AMD
EPYC audit host. ``Fresh'' and ``wrong'' are results from untouched submitted
verifier sources.}
\label{tab:results}
\begin{tabular}{lrrrrcc}
\toprule
Pair & query bytes & witness bytes & wall (s) & RSS (MiB) & fresh & wrong\\
\midrule
160S/F & 12,158  & 110 & 12.60  & 4.5  & accept & reject\\
256S/F & 31,905  & 177 & 36.44  & 10.5 & accept & reject\\
384S/F & 72,044  & 264 & 102.10 & 19.4 & accept & reject\\
512S/F & 128,850 & 352 & 8.69   & 33.7 & accept & reject\\
\bottomrule
\end{tabular}
\end{table}

A separately compiled 160-bit pipeline divided the attack into two processes.
The recovery process accepted only $(\pk,M,\sigma_S,\sigma_F)$, verified both
inputs, recovered instance zero, and checked the public RSD relation.  It took
3.99 seconds and 2,840 KiB RSS.  A second process accepted only
$(\pk,M_{\rm fresh},W)$, used fixed attacker-selected proof coins, and produced
a fresh S signature in 1.92 seconds and 4.0 MiB RSS.  The pristine verifier
accepted it.  The recovered 110-byte witness agreed with the separately retained
honest witness only in a post-recovery validation comparison.

An independent hostile implementation repeated public-only recovery and forgery
at all four levels.  As an artifact-risk control, it also compiled the untouched
submitted 160S and 160F public APIs as separate programs.  After generating one
submitted key, it deliberately reset the documented DRNG state before the two
signing calls, inducing the same nonzero $\rho$ in both untouched signers.
Those original-interface query signatures again yielded the equivalent witness,
and the fresh forgery was accepted by a separately compiled pristine verifier.
This control establishes wrapper independence only.  The untouched API normally
samples fresh $\rho$, so the control is not evidence of a vulnerability in
normal randomized use.  The hostile run's largest measured resident set across
the four levels was 46.8 MiB; this is separate from the lead-harness
measurements in Table~\ref{tab:results}.

\subsection{Real-signer sweep and negative controls}

The 256-case ReSolveD-$\alpha$-160 sweep used distinct deterministic test keys
and messages.  Every case checked honest S/F verification, IV equality, exact
witness recovery, public-relation validity, fresh-message proof generation, and
wrong-message rejection.  All 256 fresh signatures were then replayed through
the untouched submitted verifier: it accepted all 256 under their fresh
messages and rejected all 256 under their original query messages.  Per-case
wall times ranged from 12.00 to 19.36 seconds, with median 14.895 seconds; peak
RSS was at most 4,752 KiB.  These per-process times are contention-affected
rather than single-core benchmarks.

Two additional independent 160-bit cases ran the full adversarial controls.
Each control had the expected result:
\begin{itemize}[leftmargin=*,itemsep=2pt]
  \item the F opening alone could not reconstruct any of its hidden-leaf
        parents;
  \item pairing F with an S transcript from a different key, different message,
        or different $\rho$ produced a wrong witness and a rejected proof;
  \item a one-bit-mutated recovered witness produced a rejected signature;
  \item a one-bit-mutated signature was rejected by the pristine verifier; and
  \item the valid fresh forgery was rejected under the query message and an
        unrelated message.
\end{itemize}
The public recovery code has no secret-key or ground-truth input.  Secret data
is used only by the oracle processes that create the two requested signatures
and by the later equality check used to validate the experiment.

\subsection{Reproduction}

The companion repository directory contains a deterministic reproducer.  From
\texttt{reproducer/}, run
\begin{center}
\texttt{./run.sh}.
\end{center}
The command verifies the bundled source manifest, builds the conformance-repaired
query signer, public recovery program, witness-based prover, and pristine
verifier, then performs the 160-bit recovery, fresh-message forgery, and negative
controls.  A successful run ends with
\texttt{minimal\_reproducer=pass}.

The optional command \texttt{./run\_all\_levels.sh} repeats witness recovery,
fresh-message forgery, and untouched-verifier acceptance at all four levels. The
reproducer README states the compiler requirements, expected runtime and memory,
source-archive checksum, exact conformance patch, and the distinction between
the repaired signer and untouched verifier.  No audit-internal path is required.

\section{Security impact and limits}

This is a practical cross-profile composition attack enabled by absent
cryptographic separation of S/F keys and derivation domains under same-key
reuse.  Its
demonstrated consequence is recovery of an equivalent RSD signing witness
followed by a fresh-message signature accepted by the public verifier.  It is a
concrete recovery-and-forgery chain based on an equivalent signing witness,
rather than a reduced security estimate.  The original sampling seed
$\mathsf{seed}_{sk}$ is not recovered.  The recovered $W$ is nonetheless a
complete signing capability because it satisfies the same public relation
proved by an honest signer.

\paragraph{Same-key cross-profile access.}
The attacker needs S and F signatures at one level under the same full secret
key, and therefore the same effective RSD witness material.  The specification
does not instruct deployments to reuse keys across profiles.  \KeyGen{} takes a
parameter set, but the S/F profile is not cryptographically bound into or
encoded by the generated key material; both profiles accept the same serialized
key format and use the same RSD relation.  Thus the specification does not
cryptographically prevent same-key cross-profile reuse or separate the S/F key
namespaces.  A deployment that generates and enforces independent typed keys
for S and F avoids this prerequisite.

\paragraph{Same message and $\rho$.}
The two signatures must cover the same message and use the same $\rho$.  The
specified deterministic mode makes equality of $\rho$ automatic by fixing
$\rho=0^\lambda$.
The default submitted API instead normally samples fresh $\rho$; absent an
accidental equality, randomized S/F signatures do not share the root and this
attack does not apply.  The reset-DRNG experiment deliberately induced equal
nonzero $\rho$ only to test wrapper independence; it does not show a normal
randomized-API vulnerability.  Different messages likewise change $\mu$ and
the derived state.

\paragraph{No isolated EUF-CMA claim.}
The ordinary EUF-CMA game for one fixed profile provides one signing oracle.
This attack needs two profile interfaces under one key.  It therefore does not
establish an EUF-CMA break of ReSolveD-$\alpha$-160S alone, 160F alone, or any
other independently keyed single profile.

\paragraph{Opening geometry.}
The heap bounds are exact, but successful exposure of an F parent depends on the
two challenges.  Table~\ref{tab:prob} separates an exact unconditioned model from
conditioned sampling, and the 256-case signer sweep is empirical.  A failed pair
is publicly recognizable and can be retried with a new common message; the
experiments do not prove that every valid pair succeeds.

\paragraph{Specification interpretation.}
A literal reading of the PDF's Lemma~3.5 packing formula makes later QuickSilver
loop indices run outside the packed vector at 160, 256, and 384 bits and omits
230 constraint checks at 512 bits.  The submitted source instead hashes the
constraint blocks separately and does not exhibit that literal packing issue.
Xiong and Wang independently flag the printed Lemma~3.5/QS.RSDProve
block-index shift as a specification erratum~\cite{xiongwang2026}.  The
experiments use the source's natural completion.  The cross-profile leak occurs
in the preceding BAVC/VOLE layer: the recovered string independently passes the
explicit RSD relation, and fresh forgeries pass the untouched submitted
verifier.  The attack does not treat the malformed formula as a weakened proof
relation.

\section{Repair}

The direct repair is to bind the complete profile identifier into the signing
state before any tree expansion, for example
\begin{align}
 (\mu,s)&=\mathsf{XOF}(0x00\concat\mathsf{instanceID}\concat\pk\concat M;
                       3\lambda),\nonumber\\
 (r,iv)&=\mathsf{XOF}(0x01\concat\mathsf{instanceID}\concat\sk\concat\mu
                       \concat\rho;\lambda+120).               \label{eq:repair}
\end{align}
Here \textsf{instanceID} must canonically include both the security level and
the S/F profile.  The same identifier can be included in the BAVC and
Fiat--Shamir frames as defense in depth.

Keys and APIs should also be typed by the full instance.  An S key should be
rejected by the F signing entry point even though the underlying RSD relation is
the same.  Conformance tests should assert that equal raw key/message inputs in
S and F yield different roots and IVs.  Either independent key namespaces or
root domain separation blocks the demonstrated equality; deploying both makes
the intended boundary explicit.

\section{Prior art and novelty boundary}

ReSolveD introduced the RSD-specific VOLE-in-the-Head proof and its two
performance profiles~\cite{resolved-paper}.  After accounting for the live
sign-21-1 report below, a best-effort search of the NGCC inventory, the PKC
Forum, the IACR ePrint Archive, the candidate title and authors, and combinations
of ``ReSolveD'', ``cross-profile'', ``same key'', ``deterministic'', ``seed
tree'', and ``domain separation'' found no public report of the specific
cross-profile witness recovery in this paper.  This negative search is not proof
that unpublished or unindexed work does not exist.

During finalization on 1~October~2026, Xiong and Wang, with reported GLM-5.3
assistance, published a distinct one-signature multi-target attack on
ReSolveD-$\alpha$'s use of one TCCR tweak for every tree node
\cite{xiongwang2026,ngcc-resolved}.  It enumerates a hidden-parent candidate
once and tests it against the many revealed nodes of one signature, for an
expected $2^\lambda/(T+1)$ TCCR evaluations.  Median target counts give about
$2^{248.2}$, $2^{375.6}$, and $2^{503.2}$ evaluations for the affected 256-,
384-, and 512-bit sets; the 160-bit sets remain above their 128-bit target.  The
authors describe the resulting 7--9-bit reduction as non-practical: scaled
searches validate the multi-target factor and recovery chain, but the full
exponential search was not run.  The submission team nevertheless confirmed
the missing per-node tweak as a specification bug~\cite{resolved-team2026}.
That attack performs exponential same-profile enumeration.  The attack here
instead completes a hidden F leaf directly from a shared parent reconstructed
through S and runs in practical linear time.  A per-node tweak that remains
equal across S and F does not separate their root state, while an S/F
derivation tag alone does not repair within-signature TCCR target reuse; the
two fixes address different defects.

Chosen-protocol attacks give the general warning that two protocols can become
insecure when they expose a common key~\cite{kelsey1998}.  Modified-parameter
attacks similarly show that security for one parameter choice need not survive
reuse of the same key under another~\cite{howgrave2004}.  The 2024 PQC Forum
discussion of parameter-set key separation made this issue concrete for modern
signatures: Robin Larrieu suggested that even one discriminator bit between the
``s'' and ``f'' SLH-DSA profiles would suffice at a fixed hash family and
security level~\cite{larrieu2024}.  That discussion proposed preventive domain
separation; it did not give the ReSolveD-$\alpha$ attack below.

There are also close mechanistic precedents.  Bui et al. explain that reuse or
collision of a GGM master seed lets a revealed co-path complete an unopened
leaf and expose the masked witness~\cite[Sec.~3.1]{bui2024}.  Jendral and
Dubrova fault FAEST root expansion so that two VOLEitH trees coincide; when
their openings omit different leaves, one tree completes the other, permitting
reconstruction of $u$ and recovery of the witness masked by $d$
\cite[Sec.~8]{jendral2025}.  Their triggers are respectively a collision search
and a physical fault.  Jendral--Dubrova is the closest precedent for the tree
mechanism used here.  Feneuil and Rivain's formal study of correlated GGM trees
likewise makes tree correlation and domain separation explicit security
conditions~\cite{feneuil2026}.

Delgado's contemporary SDitH v2 attack is the closest standard-model and
domain-separation precedent~\cite{delgado2026}.  It also recovers a permanent
witness and forges from public cross-signature transcripts whose hidden-leaf
endpoints mask that witness.  Its mechanism and cost are different: it uses
$q=2^{12}$ ordinary same-profile signatures, exponential seed enumeration,
unary filters, and trie/Gray-code acceleration, with reported full-parameter
costs about $2^{129.957}$, $2^{193.909}$, and $2^{258.068}$; its experiments use
reduced domains.  The attack here instead needs one identical-root S/F pair and
directly completes a hidden leaf from a parent reconstructed through the
complementary opening, giving a practical linear-time recovery.  Delgado's
proposed salt, leaf, and block stream labels
do not supply ReSolveD-$\alpha$'s missing S/F derivation tag, while an S/F tag
alone would not repair Delgado's distinct expansion-domain issue.

The Galas cross-variant attack from the same NGCC audit is the direct same-key,
same-message S/F precedent~\cite{galas-attack}.  The present claim is therefore
deliberately narrow.  Its ReSolveD-specific contribution is the no-fault
trigger supplied by specified deterministic signing under unseparated S/F
namespaces, together with the shared-parent/different-leaf pcGGM identity,
completion of an arbitrary F VOLE instance through serialized corrections, and
recovery of an equivalent RSD witness rather than the original secret seed.  It
does not claim a new generic seed-tree attack family.

\section{Conclusion}

At one fixed security level, one deterministic S/F signature pair under the
same full secret key and effective RSD witness material, the same message, and
equal $\rho=0^\lambda$ can reveal an equivalent signing witness with full
signing capability.  Exposure is challenge-dependent; failure is publicly
recognizable, and the reported deterministic 160-bit sweep succeeded in
256/256 cases.  Missing profile separation makes the two executions share their
root and F internal-node seeds.  The S opening can reconstruct the parent of an
F-hidden leaf; public F bottom expansion then completes one VOLE instance,
exposes $u_0$, and unmasks the compact RSD witness.  That witness produces
ordinary fresh-message signatures.

The attack was reproduced at all four recommended levels and independently
checked through untouched submitted verifiers.  Its scope remains conditional:
it needs the same security level, full key/effective witness material, message,
and $\rho$ across S and F, with the specified deterministic mode providing equal
$\rho$.  Typed keys and full instance domain separation remove that shared state
and prevent the recovery.  This is not an isolated single-profile,
single-oracle EUF-CMA break.

\section*{Acknowledgements}

The computations were performed on the Norwegian Research and Education Cloud
(NREC), using resources provided by the University of Bergen and the University
of Oslo.

\begin{thebibliography}{99}
\small

\bibitem{resolved-spec}
Yu Yu, Kang Yang, Xiaojie Guo, Hongrui Cui, Chun Guo, Di Yan, Yituo He,
Lizheng Wang, and Jiangxia Ge,
\emph{ReSolveD-$\alpha$: Algorithm Specifications}, NGCC submission,
30 June 2026, with parameter addition dated 28 August 2026.

\bibitem{resolved-paper}
Hongrui Cui, Hanlin Liu, Di Yan, Kang Yang, Yu Yu, and Kaiyi Zhang,
\emph{ReSolveD: Shorter Signatures from Regular Syndrome Decoding and
VOLE-in-the-Head}, Public-Key Cryptography 2024; IACR Cryptology ePrint
Archive, Report 2024/040,
\url{https://eprint.iacr.org/2024/040}.

\bibitem{ngcc-resolved}
NGCC Security Reports,
\emph{ReSolveD-$\alpha$ (sign-21)},
\url{https://ngcc.dev/reports/sign-21.html}, accessed 1 October 2026.

\bibitem{xiongwang2026}
Zhenyu Xiong and Mingsheng Wang,
\emph{ReSolveD-$\alpha$---No Per-Node Tweak in the BAVC/TCCR Tree: One
Signature Gives a Multi-Target Seed Search},
PKC Forum public comment, with reported GLM-5.3 assistance, 1 October 2026,
\url{https://list.niccs.org.cn/archives/list/pkcforum@list.niccs.org.cn/message/PH6Q6UHS7YJAN3RHER7YWUNXOHOZYIFF/}.

\bibitem{resolved-team2026}
The ReSolveD-$\alpha$ Team,
\emph{Response Confirming the TCCR-Tweak Specification Bug},
PKC Forum, 1 October 2026,
\url{https://list.niccs.org.cn/archives/list/pkcforum@list.niccs.org.cn/message/GEDCFLSHBQTVSIW5K7ACXWGY7GKOXASI/}.

\bibitem{kelsey1998}
John Kelsey, Bruce Schneier, and David Wagner,
\emph{Protocol Interactions and the Chosen Protocol Attack},
Security Protocols, 5th International Workshop, Lecture Notes in Computer
Science 1361, pp.~91--104, Springer, 1998,
\url{https://doi.org/10.1007/BFb0028162}.

\bibitem{howgrave2004}
Nick Howgrave-Graham, Joseph H. Silverman, Ari Singer, and William Whyte,
\emph{Modified Parameter Attacks: Practical Attacks Against CCA2 Secure
Cryptosystems, and Countermeasures},
IACR Cryptology ePrint Archive, Report 2004/344, 2004,
\url{https://eprint.iacr.org/2004/344}.

\bibitem{larrieu2024}
Robin Larrieu,
message in \emph{Official comment on FIPS 203 ipd: seed as decapsulation key},
NIST PQC Forum, 3 June 2024,
\url{https://groups.google.com/a/list.nist.gov/g/pqc-forum/c/5CT4NC_6zRI/m/lpifFrpWAwAJ}.

\bibitem{bui2024}
Dung Bui, Eliana Carozza, Geoffroy Couteau, Dahmun Goudarzi, and Antoine Joux,
\emph{Faster Signatures from MPC-in-the-Head},
ASIACRYPT 2024; IACR Cryptology ePrint Archive, Report 2024/252,
\url{https://eprint.iacr.org/2024/252}.

\bibitem{jendral2025}
S\"onke Jendral and Elena Dubrova,
\emph{Side-Channel and Fault Injection Attacks on VOLEitH Signature
Schemes: A Case Study of Masked FAEST},
IACR Cryptology ePrint Archive, Report 2025/378, revised 2026,
\url{https://eprint.iacr.org/2025/378}.

\bibitem{feneuil2026}
Thibauld Feneuil and Matthieu Rivain,
\emph{On the Security of MPC-in-the-Head Signatures with Correlated GGM
Trees}, IACR Cryptology ePrint Archive, Report 2026/615, 2026,
\url{https://eprint.iacr.org/2026/615}.

\bibitem{delgado2026}
Jos\'e Luis Delgado,
\emph{Cross-Signature Signing-Key Recovery and Domain-Separation Repair for
SDitH v2}, IACR Cryptology ePrint Archive, Report 2026/1808, 2026,
\url{https://eprint.iacr.org/2026/1808}.

\bibitem{galas-attack}
Martin Feussner,
\emph{Cross-Variant Seed-Tree Reuse in Galas: Secret-Key Recovery and
Fresh-Message Forgery from Same-Key, Same-Message S/F Signatures},
1 October 2026,
\url{https://github.com/martinfeussner/NGCC-Signature-Audit/blob/main/Galas/Galas_Attack_Description.pdf}.

\end{thebibliography}

\end{document}
